\renewcommand{\keywords}[1]{\par\vspace{1em}{\itshape\bfseries Palabras clave -- }#1}

\chapter*{Resumen}

El mecanismo de atención constituye uno de los componentes fundamentales de las arquitecturas \textit{Transformer}, pero presenta un coste computacional cuadrático respecto a la longitud de la secuencia. FlashAttention reorganiza este cálculo para evitar materializar completamente la matriz de atención en memoria global y aprovechar de forma más eficiente la jerarquía de memoria de la GPU.

En este Trabajo Fin de Grado se estudia la implementación y optimización de FlashAttention sobre GPUs NVIDIA mediante dos niveles de abstracción diferentes. Por una parte, se desarrolla desde cero un kernel CUDA específico para arquitecturas Ampere, utilizando de forma explícita transferencias asíncronas mediante \texttt{cp.async}, cargas matriciales con \texttt{ldmatrix}, Tensor Cores mediante instrucciones \texttt{mma.sync}, memoria compartida con técnicas de \textit{swizzling} y acumulación en FP32. Por otra parte, se implementa una versión equivalente utilizando Triton, delegando al compilador una parte importante de las decisiones relacionadas con la distribución del trabajo y la gestión de memoria.

Ambas implementaciones se comparan con SDPA de PyTorch, incluyendo sus rutas FlashAttention y cuDNN. La evaluación principal se realiza sobre una NVIDIA A100-SXM4-40GB utilizando entradas BF16, dimensión de cabeza igual a 128, un único elemento por \textit{batch}, una única cabeza y atención no causal. Se utiliza un test funcional con \texttt{torch.allclose} para comprobar la corrección numérica, junto con pruebas de rendimiento para longitudes de secuencia crecientes y perfiles detallados mediante NVIDIA Nsight Compute.

Los resultados muestran que la implementación CUDA elimina los accesos globales no coalescentes y los conflictos de banco, evita \textit{register spilling} y presenta un consumo de registros y memoria compartida inferior al de varias de las rutas comparadas. Sin embargo, estas optimizaciones locales no se traducen directamente en el mejor rendimiento. El análisis revela que la granularidad del trabajo y la estrategia de paralelización tienen un impacto determinante: el kernel CUDA procesa un bloque menor de consultas por CTA, lo que reduce la reutilización de \(K\) y \(V\) y limita el paralelismo disponible para determinadas longitudes de secuencia. Triton muestra una mayor reutilización de datos. En cuDNN, el menor número de sectores solicitado es consistente con una mayor reutilización efectiva, aunque su \textit{mapping} interno no se reconstruye en este trabajo. El perfil de SDPA, por su parte, muestra una organización de ejecución distinta y un mayor número de bloques.

Para longitudes de secuencia grandes, la implementación CUDA alcanza aproximadamente \(100\,\text{TFLOP/s}\), frente a unos \(126\,\text{TFLOP/s}\) de Triton y valores superiores para SDPA y cuDNN. El trabajo muestra así que la optimización de kernels GPU no depende únicamente de minimizar accesos a memoria o consumo de recursos, sino de encontrar un equilibrio entre reutilización de datos, paralelismo, ocupación, sincronización y utilización efectiva de las unidades de cómputo.

\keywords{FlashAttention, CUDA, Triton, GPU, Tensor Cores, Ampere, optimización, Nsight Compute}
